\section{Estado Actual y Próximos Pasos}
\subsection{Lo que Ya Existe}
A la fecha de este documento (9 de octubre de 2026), las fuentes registran lo
siguiente:
\begin{itemize}
	\item El plan del proyecto se entregó el 18 de septiembre de 2026.
	\item Al 20 de septiembre estaban terminados el plan, los instrumentos de
	      medición y los avances de la investigación.
	\item El 5 de octubre de 2026, primer día del cronograma, se entregaron la
	      Ficha Inicial, el Plan de Calidad y el Tablero de Metas.
	\item El documento de investigación define el tema, el problema, la
	      pregunta, los objetivos y la metodología, con la encuesta, la entrevista y
	      la hoja de observación como apéndices.
\end{itemize}
No hay todavía mediciones del acceso, compras de materiales ni piezas del
prototipo construidas que las fuentes reporten.

\subsection{Próximos Pasos}
Según el cronograma, lo que sigue es:
\begin{enumerate}
	\item Conseguir el permiso de la dirección (16 de octubre de 2026), el
	      aviso de privacidad y la lista de materiales.
	\item Hacer el diagnóstico del 12 al 30 de octubre: encuesta, entrevistas y
	      observación del acceso durante cinco días, y compra de materiales.
	\item Pasar a las iteraciones de desarrollo a partir del 26 de octubre.
\end{enumerate}
El primer hito es el permiso de la dirección. Si no se obtiene, el plan B es
probar el prototipo en el laboratorio.

\subsection{Lo que Falta por Definir}
Estos datos no están en las fuentes y quedan pendientes:
\begin{itemize}
	\item Los resultados de todas las métricas del tablero (Tabla~\ref{tab:tablero}),
	      que dependen del diagnóstico y del piloto.
	\item La línea base del tiempo de validación, de la fila y de las incidencias
	      de identidad.
	\item El costo real, que se confirma al cotizar; el presupuesto es una
	      estimación.
	\item Los responsables de cada riesgo.
	\item Los umbrales de calidad, que son una propuesta del equipo y se ajustan
	      con la línea base del diagnóstico.
\end{itemize}
